% --- Archon physics formalization source begin ---
% archon:physics
% archon:covers PhyXMiniProblems/problem_phyx_mini_0914.lean
% archon:source-report reports/phyx_mini/problem_phyx_mini_0914.source.json

\chapter{Physics problem phyx\_mini\_0914}
\label{ch:PhyXMiniProblems_problem_phyx_mini_0914}

\paragraph{Problem source.}
\#\# Question

What is the strength of the electric field at the position indicated by the dot?

\#\# Answer choices

- A: A: 1.3 \textbackslash{}times 10\textasciicircum{}\{3\}N/C

- B: B: 5.3 \textbackslash{}times 10\textasciicircum{}\{3\}N/C

- C: C: 1.2 \textbackslash{}times 10\textasciicircum{}\{3\}N/C

- D: D: 3.5 \textbackslash{}times 10\textasciicircum{}\{3\}N/C

\#\# Figure caption (auxiliary; use the image as primary evidence)

The image shows a diagram with the following components:

1. **Two Identical Charges**: 
   - Both charges are labeled as "+3.0 nC" and are represented by red circles with a plus sign inside, indicating they are positive charges.
   - The charges are vertically aligned, one on top of the other, with a distance of "10 cm" between them.

2. **Point**:
   - There is a black dot located to the right of the upper charge.

3. **Distances**:
   - A horizontal arrow indicates a distance of "5.0 cm" from the upper charge to the black dot.
   - A vertical arrow indicates the "10 cm" separation between the two charges.

The diagram appears to illustrate an arrangement of two point charges and possibly shows a field point (black dot) where the electric field or potential might be evaluated.

\#\# Dataset metadata

Electrostatics; Physical Model Grounding Reasoning, Multi-Formula Reasoning

\paragraph{Recorded answer/context.}
C: C: 1.2 \textbackslash{}times 10\textasciicircum{}\{3\}N/C

\paragraph{Figure/image path.}
/root/proposal\_for\_physic/science-mango/phyx\_mini\_run/phyx\_data/test\_image/914.png

\paragraph{Formalization target.}
create a compiling Lean file with sorry bodies at `PhyXMiniProblems/problem\_phyx\_mini\_0914.lean`. The Lean declarations must preserve the physical quantities, dimensions or dimensional roles, figure labels, governing-law hypotheses, and final relation expressed by this problem.
Use Mathlib/Physlib names found through LeanExplore where available. If a physics API is missing, introduce faithful local abstractions rather than scalar placeholder aliases.

\begin{theorem}[Physics formalization target]
\label{thm:physics:phyx_mini_0914:target}
\lean{PhyXMiniProblems.ProblemPhyXMini0914.problem_phyx_mini_0914}
\uses{def:physics:phyx-mini-0914:phyxminiproblems-problemphyxmini0914-electricfieldstrengthinnewtonspercoulomb, def:physics:phyx-mini-0914:phyxminiproblems-problemphyxmini0914-twopointchargefieldsetup, def:physics:phyx-mini-0914:phyxminiproblems-problemphyxmini0914-matchestwostationarypointchargescenario, def:physics:phyx-mini-0914:phyxminiproblems-problemphyxmini0914-matchessuppliedtwochargefieldfigure, def:physics:phyx-mini-0914:phyxminiproblems-problemphyxmini0914-satisfiestwochargefiguregeometry, def:physics:phyx-mini-0914:phyxminiproblems-problemphyxmini0914-hasphysicaltwochargefieldparameters, def:physics:phyx-mini-0914:phyxminiproblems-problemphyxmini0914-usestextbookcoulombconstant, def:physics:phyx-mini-0914:phyxminiproblems-problemphyxmini0914-satisfiespointchargeelectricfieldlaw, def:physics:phyx-mini-0914:phyxminiproblems-problemphyxmini0914-satisfieselectricfieldsuperposition, def:physics:phyx-mini-0914:phyxminiproblems-problemphyxmini0914-calibrateselectricfieldstrengthatdot, def:physics:phyx-mini-0914:phyxminiproblems-problemphyxmini0914-roundstonearestthousand}
The assigned autoformalize agent should translate this physics problem into Lean declarations in the covered file, with theorem and lemma proof bodies written as `by sorry`.
\end{theorem}
\begin{proof}
This is an autoformalization task, not a proof task. Produce statements that can later be proved by the physics prover without weakening the physical model.
\end{proof}

% --- Archon declaration topology begin ---
\section{Declaration topology}
\begin{definition}[electric Field Strength Dimension]
  \label{def:physics:phyx-mini-0914:phyxminiproblems-problemphyxmini0914-electricfieldstrengthdimension}
  \lean{PhyXMiniProblems.ProblemPhyXMini0914.electricFieldStrengthDimension}
  The physical dimension M L T⁻² C⁻¹ of electric-field strength
\end{definition}
\begin{proof}
  This is the declared model definition of electric Field Strength Dimension.
\end{proof}
\begin{definition}[Signed Charge Quantity]
  \label{def:physics:phyx-mini-0914:phyxminiproblems-problemphyxmini0914-signedchargequantity}
  \lean{PhyXMiniProblems.ProblemPhyXMini0914.SignedChargeQuantity}
  A signed, unit-independent electric charge
\end{definition}
\begin{proof}
  This is the declared model definition of Signed Charge Quantity.
\end{proof}
\begin{definition}[Length Quantity]
  \label{def:physics:phyx-mini-0914:phyxminiproblems-problemphyxmini0914-lengthquantity}
  \lean{PhyXMiniProblems.ProblemPhyXMini0914.LengthQuantity}
  A nonnegative, unit-independent physical length
\end{definition}
\begin{proof}
  This is the declared model definition of Length Quantity.
\end{proof}
\begin{definition}[Electric Field Strength Quantity]
  \label{def:physics:phyx-mini-0914:phyxminiproblems-problemphyxmini0914-electricfieldstrengthquantity}
  \lean{PhyXMiniProblems.ProblemPhyXMini0914.ElectricFieldStrengthQuantity}
  \uses{def:physics:phyx-mini-0914:phyxminiproblems-problemphyxmini0914-electricfieldstrengthdimension}
  A nonnegative, unit-independent electric-field strength
\end{definition}
\begin{proof}
  This is the declared model definition of Electric Field Strength Quantity.
\end{proof}
\begin{definition}[charge In Coulombs]
  \label{def:physics:phyx-mini-0914:phyxminiproblems-problemphyxmini0914-chargeincoulombs}
  \lean{PhyXMiniProblems.ProblemPhyXMini0914.chargeInCoulombs}
  \uses{def:physics:phyx-mini-0914:phyxminiproblems-problemphyxmini0914-signedchargequantity}
  Read a signed physical charge in coherent-SI coulombs
\end{definition}
\begin{proof}
  This is the declared model definition of charge In Coulombs.
\end{proof}
\begin{definition}[charge In Nanocoulombs]
  \label{def:physics:phyx-mini-0914:phyxminiproblems-problemphyxmini0914-chargeinnanocoulombs}
  \lean{PhyXMiniProblems.ProblemPhyXMini0914.chargeInNanocoulombs}
  \uses{def:physics:phyx-mini-0914:phyxminiproblems-problemphyxmini0914-signedchargequantity, def:physics:phyx-mini-0914:phyxminiproblems-problemphyxmini0914-chargeincoulombs}
  Read a charge in the nanocoulombs printed in the figure
\end{definition}
\begin{proof}
  This is the declared model definition of charge In Nanocoulombs.
\end{proof}
\begin{definition}[length In Meters]
  \label{def:physics:phyx-mini-0914:phyxminiproblems-problemphyxmini0914-lengthinmeters}
  \lean{PhyXMiniProblems.ProblemPhyXMini0914.lengthInMeters}
  \uses{def:physics:phyx-mini-0914:phyxminiproblems-problemphyxmini0914-lengthquantity}
  Read a nonnegative physical length in coherent-SI metres
\end{definition}
\begin{proof}
  This is the declared model definition of length In Meters.
\end{proof}
\begin{definition}[length In Centimeters]
  \label{def:physics:phyx-mini-0914:phyxminiproblems-problemphyxmini0914-lengthincentimeters}
  \lean{PhyXMiniProblems.ProblemPhyXMini0914.lengthInCentimeters}
  \uses{def:physics:phyx-mini-0914:phyxminiproblems-problemphyxmini0914-lengthquantity, def:physics:phyx-mini-0914:phyxminiproblems-problemphyxmini0914-lengthinmeters}
  Read a physical length in the centimetres used by the two arrows
\end{definition}
\begin{proof}
  This is the declared model definition of length In Centimeters.
\end{proof}
\begin{definition}[electric Field Strength In Newtons Per Coulomb]
  \label{def:physics:phyx-mini-0914:phyxminiproblems-problemphyxmini0914-electricfieldstrengthinnewtonspercoulomb}
  \lean{PhyXMiniProblems.ProblemPhyXMini0914.electricFieldStrengthInNewtonsPerCoulomb}
  \uses{def:physics:phyx-mini-0914:phyxminiproblems-problemphyxmini0914-electricfieldstrengthquantity}
  Read a field strength in coherent-SI newtons per coulomb
\end{definition}
\begin{proof}
  This is the declared model definition of electric Field Strength In Newtons Per Coulomb.
\end{proof}
\begin{definition}[Source Charge]
  \label{def:physics:phyx-mini-0914:phyxminiproblems-problemphyxmini0914-sourcecharge}
  \lean{PhyXMiniProblems.ProblemPhyXMini0914.SourceCharge}
  The two identical source charges visible in the primary image
\end{definition}
\begin{proof}
  This is the declared model definition of Source Charge.
\end{proof}
\begin{definition}[Diagram Site]
  \label{def:physics:phyx-mini-0914:phyxminiproblems-problemphyxmini0914-diagramsite}
  \lean{PhyXMiniProblems.ProblemPhyXMini0914.DiagramSite}
  The two source locations and the black observation dot
\end{definition}
\begin{proof}
  This is the declared model definition of Diagram Site.
\end{proof}
\begin{definition}[to Diagram Site]
  \label{def:physics:phyx-mini-0914:phyxminiproblems-problemphyxmini0914-sourcecharge-todiagramsite}
  \lean{PhyXMiniProblems.ProblemPhyXMini0914.SourceCharge.toDiagramSite}
  \uses{def:physics:phyx-mini-0914:phyxminiproblems-problemphyxmini0914-sourcecharge, def:physics:phyx-mini-0914:phyxminiproblems-problemphyxmini0914-diagramsite}
  Regard a source charge as its corresponding geometric site
\end{definition}
\begin{proof}
  This is the declared model definition of to Diagram Site.
\end{proof}
\begin{definition}[Figure Charge Sign]
  \label{def:physics:phyx-mini-0914:phyxminiproblems-problemphyxmini0914-figurechargesign}
  \lean{PhyXMiniProblems.ProblemPhyXMini0914.FigureChargeSign}
  The sign glyph drawn inside each red charge marker
\end{definition}
\begin{proof}
  This is the declared model definition of Figure Charge Sign.
\end{proof}
\begin{definition}[Two Charge Field Figure]
  \label{def:physics:phyx-mini-0914:phyxminiproblems-problemphyxmini0914-twochargefieldfigure}
  \lean{PhyXMiniProblems.ProblemPhyXMini0914.TwoChargeFieldFigure}
  \uses{def:physics:phyx-mini-0914:phyxminiproblems-problemphyxmini0914-sourcecharge, def:physics:phyx-mini-0914:phyxminiproblems-problemphyxmini0914-figurechargesign}
  Literal presentation data transcribed from image 914.png
\end{definition}
\begin{proof}
  This is the declared model definition of Two Charge Field Figure.
\end{proof}
\begin{definition}[Two Point Charge Field Setup]
  \label{def:physics:phyx-mini-0914:phyxminiproblems-problemphyxmini0914-twopointchargefieldsetup}
  \lean{PhyXMiniProblems.ProblemPhyXMini0914.TwoPointChargeFieldSetup}
  \uses{def:physics:phyx-mini-0914:phyxminiproblems-problemphyxmini0914-signedchargequantity, def:physics:phyx-mini-0914:phyxminiproblems-problemphyxmini0914-lengthquantity, def:physics:phyx-mini-0914:phyxminiproblems-problemphyxmini0914-electricfieldstrengthquantity, def:physics:phyx-mini-0914:phyxminiproblems-problemphyxmini0914-sourcecharge, def:physics:phyx-mini-0914:phyxminiproblems-problemphyxmini0914-diagramsite, def:physics:phyx-mini-0914:phyxminiproblems-problemphyxmini0914-twochargefieldfigure}
  Independent physical quantities and observables for the two-charge system. The planar coordinates are measured in metres in an arbitrary orthonormal frame whose rightward and upward directions are also recorded explicitly
\end{definition}
\begin{proof}
  This is the declared model definition of Two Point Charge Field Setup.
\end{proof}
\begin{definition}[displacement From Source To Dot In Meters]
  \label{def:physics:phyx-mini-0914:phyxminiproblems-problemphyxmini0914-displacementfromsourcetodotinmeters}
  \lean{PhyXMiniProblems.ProblemPhyXMini0914.displacementFromSourceToDotInMeters}
  \uses{def:physics:phyx-mini-0914:phyxminiproblems-problemphyxmini0914-sourcecharge, def:physics:phyx-mini-0914:phyxminiproblems-problemphyxmini0914-twopointchargefieldsetup}
  Displacement in metres from a source charge to the observation dot
\end{definition}
\begin{proof}
  This is the declared model definition of displacement From Source To Dot In Meters.
\end{proof}
\begin{definition}[source Field Vector At Dot]
  \label{def:physics:phyx-mini-0914:phyxminiproblems-problemphyxmini0914-sourcefieldvectoratdot}
  \lean{PhyXMiniProblems.ProblemPhyXMini0914.sourceFieldVectorAtDot}
  \uses{def:physics:phyx-mini-0914:phyxminiproblems-problemphyxmini0914-sourcecharge, def:physics:phyx-mini-0914:phyxminiproblems-problemphyxmini0914-twopointchargefieldsetup}
  The source contribution evaluated at the marked dot and observation time
\end{definition}
\begin{proof}
  This is the declared model definition of source Field Vector At Dot.
\end{proof}
\begin{definition}[total Field Vector At Dot]
  \label{def:physics:phyx-mini-0914:phyxminiproblems-problemphyxmini0914-totalfieldvectoratdot}
  \lean{PhyXMiniProblems.ProblemPhyXMini0914.totalFieldVectorAtDot}
  \uses{def:physics:phyx-mini-0914:phyxminiproblems-problemphyxmini0914-twopointchargefieldsetup}
  The total electric-field vector evaluated at the marked dot
\end{definition}
\begin{proof}
  This is the declared model definition of total Field Vector At Dot.
\end{proof}
\begin{definition}[Matches Two Stationary Point Charge Scenario]
  \label{def:physics:phyx-mini-0914:phyxminiproblems-problemphyxmini0914-matchestwostationarypointchargescenario}
  \lean{PhyXMiniProblems.ProblemPhyXMini0914.MatchesTwoStationaryPointChargeScenario}
  \uses{def:physics:phyx-mini-0914:phyxminiproblems-problemphyxmini0914-twopointchargefieldsetup}
  The two pictured sources are stationary point charges
\end{definition}
\begin{proof}
  This is the declared model definition of Matches Two Stationary Point Charge Scenario.
\end{proof}
\begin{definition}[Matches Supplied Two Charge Field Figure]
  \label{def:physics:phyx-mini-0914:phyxminiproblems-problemphyxmini0914-matchessuppliedtwochargefieldfigure}
  \lean{PhyXMiniProblems.ProblemPhyXMini0914.MatchesSuppliedTwoChargeFieldFigure}
  \uses{def:physics:phyx-mini-0914:phyxminiproblems-problemphyxmini0914-chargeinnanocoulombs, def:physics:phyx-mini-0914:phyxminiproblems-problemphyxmini0914-lengthincentimeters, def:physics:phyx-mini-0914:phyxminiproblems-problemphyxmini0914-twopointchargefieldsetup}
  Every qualitative feature and numerical label read directly from the primary image, together with the calibration of those labels to physical quantities. No electric-field value occurs in this predicate
\end{definition}
\begin{proof}
  This is the declared model definition of Matches Supplied Two Charge Field Figure.
\end{proof}
\begin{definition}[Satisfies Two Charge Figure Geometry]
  \label{def:physics:phyx-mini-0914:phyxminiproblems-problemphyxmini0914-satisfiestwochargefiguregeometry}
  \lean{PhyXMiniProblems.ProblemPhyXMini0914.SatisfiesTwoChargeFigureGeometry}
  \uses{def:physics:phyx-mini-0914:phyxminiproblems-problemphyxmini0914-lengthinmeters, def:physics:phyx-mini-0914:phyxminiproblems-problemphyxmini0914-twopointchargefieldsetup}
  The metric geometry depicted by the image. The observation dot is displaced rightward from the upper charge, and the upper charge is displaced upward from the lower charge. The two indicated directions form an orthonormal frame
\end{definition}
\begin{proof}
  This is the declared model definition of Satisfies Two Charge Figure Geometry.
\end{proof}
\begin{definition}[Has Physical Two Charge Field Parameters]
  \label{def:physics:phyx-mini-0914:phyxminiproblems-problemphyxmini0914-hasphysicaltwochargefieldparameters}
  \lean{PhyXMiniProblems.ProblemPhyXMini0914.HasPhysicalTwoChargeFieldParameters}
  \uses{def:physics:phyx-mini-0914:phyxminiproblems-problemphyxmini0914-chargeincoulombs, def:physics:phyx-mini-0914:phyxminiproblems-problemphyxmini0914-lengthinmeters, def:physics:phyx-mini-0914:phyxminiproblems-problemphyxmini0914-twopointchargefieldsetup, def:physics:phyx-mini-0914:phyxminiproblems-problemphyxmini0914-displacementfromsourcetodotinmeters}
  Positivity and nondegeneracy conditions for the electrostatic model
\end{definition}
\begin{proof}
  This is the declared model definition of Has Physical Two Charge Field Parameters.
\end{proof}
\begin{definition}[Uses Textbook Coulomb Constant]
  \label{def:physics:phyx-mini-0914:phyxminiproblems-problemphyxmini0914-usestextbookcoulombconstant}
  \lean{PhyXMiniProblems.ProblemPhyXMini0914.UsesTextbookCoulombConstant}
  \uses{def:physics:phyx-mini-0914:phyxminiproblems-problemphyxmini0914-twopointchargefieldsetup}
  The rounded value of Coulomb's constant customarily used for this textbook multiple-choice calculation. It is a background constant, not a field-value answer
\end{definition}
\begin{proof}
  This is the declared model definition of Uses Textbook Coulomb Constant.
\end{proof}
\begin{definition}[Satisfies Point Charge Electric Field Law]
  \label{def:physics:phyx-mini-0914:phyxminiproblems-problemphyxmini0914-satisfiespointchargeelectricfieldlaw}
  \lean{PhyXMiniProblems.ProblemPhyXMini0914.SatisfiesPointChargeElectricFieldLaw}
  \uses{def:physics:phyx-mini-0914:phyxminiproblems-problemphyxmini0914-chargeincoulombs, def:physics:phyx-mini-0914:phyxminiproblems-problemphyxmini0914-twopointchargefieldsetup, def:physics:phyx-mini-0914:phyxminiproblems-problemphyxmini0914-displacementfromsourcetodotinmeters, def:physics:phyx-mini-0914:phyxminiproblems-problemphyxmini0914-sourcefieldvectoratdot}
  For a source charge q and source-to-dot displacement r, the contribution at the dot is k q r / ‖r‖³. This vector equation combines the inverse-square magnitude with the radial direction and contains no answer-choice value
\end{definition}
\begin{proof}
  This is the declared model definition of Satisfies Point Charge Electric Field Law.
\end{proof}
\begin{definition}[Satisfies Electric Field Superposition]
  \label{def:physics:phyx-mini-0914:phyxminiproblems-problemphyxmini0914-satisfieselectricfieldsuperposition}
  \lean{PhyXMiniProblems.ProblemPhyXMini0914.SatisfiesElectricFieldSuperposition}
  \uses{def:physics:phyx-mini-0914:phyxminiproblems-problemphyxmini0914-twopointchargefieldsetup, def:physics:phyx-mini-0914:phyxminiproblems-problemphyxmini0914-sourcefieldvectoratdot, def:physics:phyx-mini-0914:phyxminiproblems-problemphyxmini0914-totalfieldvectoratdot}
  The total field at the dot is the vector sum of the two source fields
\end{definition}
\begin{proof}
  This is the declared model definition of Satisfies Electric Field Superposition.
\end{proof}
\begin{definition}[Calibrates Electric Field Strength At Dot]
  \label{def:physics:phyx-mini-0914:phyxminiproblems-problemphyxmini0914-calibrateselectricfieldstrengthatdot}
  \lean{PhyXMiniProblems.ProblemPhyXMini0914.CalibratesElectricFieldStrengthAtDot}
  \uses{def:physics:phyx-mini-0914:phyxminiproblems-problemphyxmini0914-electricfieldstrengthinnewtonspercoulomb, def:physics:phyx-mini-0914:phyxminiproblems-problemphyxmini0914-twopointchargefieldsetup, def:physics:phyx-mini-0914:phyxminiproblems-problemphyxmini0914-totalfieldvectoratdot}
  The dimensionful scalar observable is calibrated to the Euclidean norm of the total Physlib electric-field vector at the dot
\end{definition}
\begin{proof}
  This is the declared model definition of Calibrates Electric Field Strength At Dot.
\end{proof}
\begin{lemma}[upper Source Distance In Meters]
  \label{lem:physics:phyx-mini-0914:phyxminiproblems-problemphyxmini0914-uppersourcedistanceinmeters}
  \lean{PhyXMiniProblems.ProblemPhyXMini0914.upperSourceDistanceInMeters}
  \uses{def:physics:phyx-mini-0914:phyxminiproblems-problemphyxmini0914-twopointchargefieldsetup, def:physics:phyx-mini-0914:phyxminiproblems-problemphyxmini0914-displacementfromsourcetodotinmeters, def:physics:phyx-mini-0914:phyxminiproblems-problemphyxmini0914-matchessuppliedtwochargefieldfigure, def:physics:phyx-mini-0914:phyxminiproblems-problemphyxmini0914-satisfiestwochargefiguregeometry}
  The upper source is the displayed 5.0 cm from the dot
\end{lemma}
\begin{proof}
  The conclusion follows from the governing relations and figure readouts named in the statement of upper Source Distance In Meters.
\end{proof}
\begin{lemma}[lower Source Displacement In Meters]
  \label{lem:physics:phyx-mini-0914:phyxminiproblems-problemphyxmini0914-lowersourcedisplacementinmeters}
  \lean{PhyXMiniProblems.ProblemPhyXMini0914.lowerSourceDisplacementInMeters}
  \uses{def:physics:phyx-mini-0914:phyxminiproblems-problemphyxmini0914-lengthinmeters, def:physics:phyx-mini-0914:phyxminiproblems-problemphyxmini0914-twopointchargefieldsetup, def:physics:phyx-mini-0914:phyxminiproblems-problemphyxmini0914-displacementfromsourcetodotinmeters, def:physics:phyx-mini-0914:phyxminiproblems-problemphyxmini0914-satisfiestwochargefiguregeometry}
  The lower-source displacement is the sum of the displayed rightward 5.0 cm offset and upward 10 cm separation
\end{lemma}
\begin{proof}
  The conclusion follows from the governing relations and figure readouts named in the statement of lower Source Displacement In Meters.
\end{proof}
\begin{definition}[Answer Choice]
  \label{def:physics:phyx-mini-0914:phyxminiproblems-problemphyxmini0914-answerchoice}
  \lean{PhyXMiniProblems.ProblemPhyXMini0914.AnswerChoice}
  Labels attached to the four displayed field-strength choices
\end{definition}
\begin{proof}
  This is the declared model definition of Answer Choice.
\end{proof}
\begin{definition}[field Strength In Newtons Per Coulomb]
  \label{def:physics:phyx-mini-0914:phyxminiproblems-problemphyxmini0914-answerchoice-fieldstrengthinnewtonspercoulomb}
  \lean{PhyXMiniProblems.ProblemPhyXMini0914.AnswerChoice.fieldStrengthInNewtonsPerCoulomb}
  \uses{def:physics:phyx-mini-0914:phyxminiproblems-problemphyxmini0914-answerchoice}
  Field strength in newtons per coulomb printed beside each answer label
\end{definition}
\begin{proof}
  This is the declared model definition of field Strength In Newtons Per Coulomb.
\end{proof}
\begin{definition}[recorded Dataset Answer]
  \label{def:physics:phyx-mini-0914:phyxminiproblems-problemphyxmini0914-recordeddatasetanswer}
  \lean{PhyXMiniProblems.ProblemPhyXMini0914.recordedDatasetAnswer}
  \uses{def:physics:phyx-mini-0914:phyxminiproblems-problemphyxmini0914-answerchoice}
  The answer label recorded by the supplied dataset metadata
\end{definition}
\begin{proof}
  This is the declared model definition of recorded Dataset Answer.
\end{proof}
\begin{definition}[Rounds To Nearest Thousand]
  \label{def:physics:phyx-mini-0914:phyxminiproblems-problemphyxmini0914-roundstonearestthousand}
  \lean{PhyXMiniProblems.ProblemPhyXMini0914.RoundsToNearestThousand}
  value rounds to rounded to the nearest thousand. Dividing by 1000 and using Mathlib's round makes the intended two-significant-figure convention for a value near 1.2 * 10\textasciicircum{}4 explicit
\end{definition}
\begin{proof}
  This is the declared model definition of Rounds To Nearest Thousand.
\end{proof}
% --- Archon declaration topology end ---
% --- Archon physics formalization source end ---
